\chapter{Saving and resuming a session}
\label{ch:session}

On a big design, loading takes time, and you may want to try a new script again and
again after the same first two. A \term{session} is the state of a run: the text of every
file with the edits so far, the expectations and checks of the operations, the log, the
notes, the added files and the scripts that ran. You can save it and continue later,
without the input tree and without running the first scripts again.

\section{Saving}

\opt{--save} \code{FILE} writes the session to a zip archive after all scripts have run. The
archive holds plain data, and nothing in it is code: loading it never runs anything.
You can open it with any zip tool. It has \file{session.json} (the index),
\file{original/} (every source file as it was loaded), \file{current/} (the files that
differ) and \file{added/} (the files that scripts added). The same state gives the same
bytes. An existing archive is replaced.

\begin{itemize}
\item Without \opt{-o}, hdl-vet saves, prints \code{saved FILE} and the summary, and then
  one line, \code{sign-off so far: ok} or \code{sign-off so far: FAILED}. After
  \code{FAILED} the report follows on the error output. The exit status is 0 in both
  cases: saving a state that does not pass is not an error. The state is saved as it is.
\item With \opt{-o} or \opt{--dry-run} too, the archive is written once every script has
  run, before the sign-off. The run then goes on as usual: it can still exit with
  status 1, and the archive exists either way. The line \code{sign-off so far} is not
  printed.
\item A script that fails writes no archive.
\item hdl-vet refuses a \opt{--save} file that lies inside the input tree or inside the
  \opt{-o} folder, or that is a script, the project file or a file list that the design
  reads: the top-level one, or one that it reads with \code{-f} or \code{-F}, wherever
  that list lies. A resumed session refuses the same file lists.
  The scripts are then not run: \code{apply} is not called. The code at the top level of
  every script has already run by then, because hdl-vet reads each script before it
  checks \opt{--save}.
\end{itemize}

\section{Resuming}

\opt{--resume} \code{FILE} continues from a saved session instead of reading a file list. Do
not give \opt{-f}, \opt{--top} or \opt{--base}: the session has them, and hdl-vet refuses
them. You add scripts with \opt{-s}, and you may save again with \opt{--save}; that
replaces the old archive when it names the same file.

\begin{itemize}
\item A script that the session has applied already (same name, same content) is
  refused.
\item The summary counts only the operations since the resume, and only the notes since
  the resume are printed.
\item The sign-off covers every script, the old ones and the new. The header line and the
  manifest list all scripts, the saved ones first. If you save the state and resume it
  with one more script, the output tree is byte for byte the tree of one run of all the
  scripts.
\item The session reports and lays out files exactly as the first run did. It needs the
  same folders: it records where the input tree was.
\end{itemize}

With a project file, the project's scripts are compared with the saved list from the
start, by name and content. If the project's script number $n$ is not the saved script
number $n$, hdl-vet stops and says that the project has changed since the save.
\opt{--no-project} resumes regardless. Project scripts that the session does not have
yet run first, then the scripts of \opt{-s}.

\lstinputlisting[style=console]{examples/console/session.txt}

\section{What is and is not detected}

A saved session is used as it is. hdl-vet does not check whether the file list, the
sources or the scripts changed after the save. In the transcript, the command \cmd{cp}
puts a regenerated \file{soc_top.sv} into the design after the save. The resume does not
notice. It continues from the text that it saved, so the output has no \code{busy_o}.
If your generator ran again, start from the beginning.

What hdl-vet does detect:

\begin{itemize}
\item A script that has been applied already (above).
\item An archive that is not a session archive, is damaged, or has another format
  version. The message names the file.
\item A project whose scripts differ from the saved ones (above).
\item \opt{-f}, \opt{--top} or \opt{--base} together with \opt{--resume}.
\end{itemize}

A scope or a group that is open cannot be saved (chapter~\ref{ch:scripts}). Nothing about
scopes is kept in a session.
